% =============================================================================
%  Section 3 — Direct Repository Access (Raw Runs)
%  Drop into the full report between Section 2 (MCP) and Section 4 (Comparison)
% =============================================================================

\section{Direct Repository Access Evaluation (Raw Runs)}\label{sec:raw}

\subsection{Methodology}

Each of the 22 \texttt{astropy/astropy} tasks from the
\texttt{princeton-nlp/SWE-bench\_Verified} dataset was presented to
Claude Code (model: \texttt{claude-opus-4-6}) with \emph{direct
filesystem access} to a locally cloned copy of the repository, checked
out at the exact base commit specified by the task.
No MCP knowledge graph, no web search, no \texttt{git} CLI, and no
pre-computed answer cache were available to the model.

\paragraph{Prompt structure.}
Every task was issued under a fixed set of ground rules:

\begin{enumerate}
  \item \textbf{Read files directly} from the local clone at the base-commit path.
  \item \textbf{Form the answer entirely from first-principles evidence}
        gathered during the session.
  \item \textbf{No web search}, \textbf{no \texttt{git} commands}
        (\texttt{git log}, \texttt{git blame}, \texttt{git diff}, etc.),
        and \textbf{no remote API calls}.
  \item \textbf{No reading} from \texttt{answers/}, \texttt{results/},
        \texttt{expected/}, or \texttt{cached/} directories.
\end{enumerate}

\noindent
The SWE-Bench question text (issue description, hints, and required
failing tests) was appended after the ground rules.
Claude Code was then free to explore the codebase, read any source file,
and produce a unified diff patch as its answer, which was written to
\texttt{answer.json} alongside telemetry (time, cost, per-model token
counts).

\paragraph{Judging.}
Answers were graded against the ground-truth patch on a
three-dimensional rubric (Table~\ref{tab:rubric-raw}), scored by a
separate Claude Opus~4.6 judge session.

\begin{table}[h]
\centering
\caption{Scoring rubric — 8-point scale.}\label{tab:rubric-raw}
\begin{tabular}{lcp{8cm}}
\toprule
\textbf{Dimension} & \textbf{Max} & \textbf{Criterion} \\
\midrule
Root cause identification & 3 & Correctly diagnoses \emph{why} the bug exists. \\
Correct file(s)           & 2 & Identifies the right file(s) to change. \\
Correct patch / code change & 3 & Proposed change matches ground truth or is functionally equivalent. \\
\bottomrule
\end{tabular}
\end{table}

\noindent
Grade tiers: \textbf{Exact}~(8/8), \textbf{Near-perfect}~(7/8),
\textbf{Partial}~(5--6/8), \textbf{Fail}~($\le$4/8).

% ---------------------------------------------------------------------------
\subsection{Results}

\subsubsection{Aggregate Performance}

\begin{table}[h]
\centering
\caption{Dimension-level scores across all 22 tasks.}\label{tab:dim-raw}
\begin{tabular}{lrrr}
\toprule
\textbf{Dimension} & \textbf{Score} & \textbf{Max} & \textbf{\%} \\
\midrule
Root cause identification   & 64 & 66 & 97.0\% \\
Correct file(s)             & 44 & 44 & 100\%  \\
Correct patch / code change & 54 & 66 & 81.8\% \\
\midrule
\textbf{Overall}            & \textbf{162} & \textbf{176} & \textbf{92.0\%} \\
\bottomrule
\end{tabular}
\end{table}

\begin{table}[h]
\centering
\caption{Grade distribution.}\label{tab:grade-raw}
\begin{tabular}{lcl}
\toprule
\textbf{Grade} & \textbf{Count} & \textbf{Tasks} \\
\midrule
Exact (8/8)        & 14 & 1--6, 9, 12--18 \\
Near-perfect (7/8) &  5 & 8, 10, 11, 20, 22 \\
Partial (5--6/8)   &  2 & 7, 21 \\
Fail ($\le$4/8)    &  1 & 19 \\
\bottomrule
\end{tabular}
\end{table}

\subsubsection{Per-Task Scores, Time \& Cost}

\begin{table}[h]
\centering
\small
\caption{Per-task results — score, wall-clock time, and API cost.}\label{tab:per-task-raw}
\setlength{\tabcolsep}{4pt}
\begin{tabular}{clccccrr}
\toprule
\textbf{\#} & \textbf{Instance ID} & \textbf{RC} & \textbf{Files} & \textbf{Patch} & \textbf{Score} & \textbf{Time (s)} & \textbf{Cost (\$)} \\
\midrule
 1 & \texttt{astropy-12907} & 3 & 2 & 3 & 8/8 &   259 & 0.400 \\
 2 & \texttt{astropy-13033} & 3 & 2 & 3 & 8/8 &   191 & 0.669 \\
 3 & \texttt{astropy-13236} & 3 & 2 & 3 & 8/8 &   422 & 1.395 \\
 4 & \texttt{astropy-13398} & 3 & 2 & 3 & 8/8 &   612 & 1.772 \\
 5 & \texttt{astropy-13453} & 3 & 2 & 3 & 8/8 &   123 & 0.635 \\
 6 & \texttt{astropy-13579} & 3 & 2 & 3 & 8/8 &   223 & 0.741 \\
 7 & \texttt{astropy-13977} & 3 & 2 & 1 & 6/8 &   212 & 0.787 \\
 8 & \texttt{astropy-14096} & 3 & 2 & 2 & 7/8 &   175 & 0.619 \\
 9 & \texttt{astropy-14182} & 3 & 2 & 3 & 8/8 &   232 & 0.565 \\
10 & \texttt{astropy-14309} & 3 & 2 & 2 & 7/8 &   171 & 0.153 \\
11 & \texttt{astropy-14365} & 3 & 2 & 2 & 7/8 &   141 & 0.496 \\
12 & \texttt{astropy-14369} & 3 & 2 & 3 & 8/8 &   475 & 1.492 \\
13 & \texttt{astropy-14508} & 3 & 2 & 3 & 8/8 &   100 & 0.385 \\
14 & \texttt{astropy-14539} & 3 & 2 & 3 & 8/8 &   123 & 0.511 \\
15 & \texttt{astropy-14598} & 3 & 2 & 3 & 8/8 &   600 & 1.768 \\
16 & \texttt{astropy-14995} & 3 & 2 & 3 & 8/8 &   153 & 0.510 \\
17 & \texttt{astropy-7166}  & 3 & 2 & 3 & 8/8 &   151 & 0.393 \\
18 & \texttt{astropy-7336}  & 3 & 2 & 3 & 8/8 &   291 & 0.307 \\
19 & \texttt{astropy-7606}  & 1 & 2 & 0 & 3/8 &   181 & 0.392 \\
20 & \texttt{astropy-7671}  & 3 & 2 & 2 & 7/8 &   202 & 0.682 \\
21 & \texttt{astropy-8707}  & 3 & 2 & 1 & 6/8 &   372 & 1.135 \\
22 & \texttt{astropy-8872}  & 3 & 2 & 2 & 7/8 &   119 & 0.353 \\
\midrule
   & \textbf{Total}  & \textbf{64/66} & \textbf{44/44} & \textbf{54/66} & \textbf{162/176} & \textbf{5{,}528} & \textbf{16.16} \\
   & \textbf{Average} &       &       &       & \textbf{7.4/8}   & \textbf{251} & \textbf{0.73} \\
\bottomrule
\end{tabular}
\end{table}

% ---------------------------------------------------------------------------
\subsubsection{Token \& Cost Breakdown (Aggregated)}

\begin{table}[h]
\centering
\small
\caption{Aggregated token usage and cost across all 22 tasks.}\label{tab:tokens-agg-raw}
\setlength{\tabcolsep}{4pt}
\begin{tabular}{lrrrrrrr}
\toprule
& \textbf{Input} & \textbf{Cache Write} & \textbf{Cache Read} & \textbf{Output}
& \textbf{Eff Weighted} & \textbf{Eff Input} & \textbf{Cost (\$)} \\
\midrule
Total   & 65{,}038 & 1{,}057{,}447 & 14{,}794{,}821 & 276{,}939 & 2{,}866{,}329 & 15{,}917{,}306 & 16.16 \\
Average & 2{,}956  & 48{,}066      & 672{,}492       & 12{,}588  & 130{,}288     & 723{,}514      & 0.73  \\
\bottomrule
\end{tabular}
\\[4pt]
{\footnotesize
Effective Weighted Input = Input + $1.25 \times$ Cache Write + $0.1 \times$ Cache Read.\quad
Effective Input = Input + Cache Write + Cache Read.}
\end{table}

% ---------------------------------------------------------------------------
\subsubsection{Per-Model Token Usage}

Claude Code internally dispatches API requests to two models:
\texttt{claude-opus-4-6} for reasoning and patch generation, and
\texttt{claude-haiku-4-5} for lightweight triage and tool routing.
Table~\ref{tab:model-split-raw} shows the split.

\begin{table}[h]
\centering
\caption{Token and cost split between Opus and Haiku.}\label{tab:model-split-raw}
\begin{tabular}{lrrrr}
\toprule
\textbf{Metric} & \textbf{Opus} & \textbf{Haiku} & \textbf{Total} & \textbf{Opus \%} \\
\midrule
Input tokens       &          506 &     64{,}532 &      65{,}038 &   0.8\% \\
Cache Write tokens &    642{,}844 &    414{,}603 &   1{,}057{,}447 &  60.8\% \\
Cache Read tokens  & 10{,}483{,}061 & 4{,}311{,}760 & 14{,}794{,}821 &  70.9\% \\
Output tokens      &    224{,}986 &     51{,}953 &     276{,}939 &  81.2\% \\
Cost (USD)         &     \$14.886 &      \$1.274 &      \$16.160 &  92.1\% \\
API requests       &         350 &          178 &          528 &  66.3\% \\
\bottomrule
\end{tabular}
\end{table}

% ---------------------------------------------------------------------------
\subsection{Observations}\label{sec:raw-obs}

The following observations are derived strictly from the results
presented above.

\begin{enumerate}

\item \textbf{High overall accuracy.}
      The model scored 162/176 (92.0\%), with 14 of 22 tasks receiving
      a perfect 8/8 and only a single outright failure (task~19,
      \texttt{astropy-7606}, scoring~3/8).

\item \textbf{Root cause identification is near-flawless.}
      The root-cause dimension scored 64/66 (97.0\%).  The only
      deduction came from task~19 where the model returned a filename
      but no diagnosis.  This indicates that reading the full codebase
      is sufficient for Claude Opus~4.6 to reliably pinpoint the
      source of almost every bug.

\item \textbf{File identification is perfect.}
      All 22 tasks received the maximum file-identification score
      (44/44).  Direct repository access eliminates any ambiguity about
      which files exist and where they reside.

\item \textbf{Patch quality is the limiting dimension.}
      The patch dimension scored 54/66 (81.8\%), accounting for all
      lost points except those from task~19.  Common deduction patterns
      include: incomplete scope of the fix (task~7 wrapping only part
      of the function body), description-only answers without an actual
      diff (task~21), and minor semantic divergences from the
      ground-truth approach that are functionally equivalent but less
      precise (tasks~8, 10, 11, 20, 22).

\item \textbf{Opus dominates cost; Haiku acts as a triage layer.}
      Opus accounts for 92.1\% of total cost (\$14.89 of \$16.16) and
      generates 81.2\% of all output tokens, yet receives only 0.8\%
      of raw input tokens.  Haiku absorbs 99.2\% of input tokens and
      appears to serve as a lightweight routing or summarisation layer,
      invoking Opus for the substantive reasoning steps.  On 11 of 22
      tasks, Haiku made exactly one API request with zero cache
      activity, suggesting a single triage call before handing off
      entirely to Opus.

\item \textbf{Cost is moderate and predictable.}
      Average cost per task is \$0.73 (median lower, range
      \$0.15--\$1.77).  The two costliest tasks (\texttt{13398} at
      \$1.77 and \texttt{14598} at \$1.77) are also among the longest
      running, correlating cost with codebase exploration depth rather
      than task difficulty rating.

\item \textbf{Wall-clock time averages 251\,s per task.}
      Total evaluation time was 5{,}528\,s (\(\approx\)92\,min) for
      22~tasks.  The fastest task (\texttt{14508}) completed in 100\,s;
      the slowest (\texttt{13398}) took 612\,s.  Tasks rated
      ``1--4\,h'' human difficulty did not consistently take longer
      than ``15\,min--1\,h'' tasks, suggesting that model difficulty
      does not map linearly to human difficulty estimates.

\item \textbf{Cache read tokens dominate the token budget.}
      Cache reads constitute 91.4\% of all effective input tokens
      (14.8\,M of 15.9\,M), indicating heavy re-reading of previously
      loaded context as the model iterates through exploration and
      patch generation cycles.

\item \textbf{The sole failure (task~19) produced no patch.}
      Task~19 (\texttt{astropy-7606}) required returning
      \texttt{NotImplemented} from two locations in
      \texttt{units/core.py}.  The model identified the correct file
      but produced only a filename with no root-cause analysis and no
      diff, scoring 3/8.  This is the only task where the model failed
      to attempt a code change.

\item \textbf{Re-runs improved two tasks from 7/8 to 8/8.}
      Tasks~4 (\texttt{13398}) and~5 (\texttt{13453}) were re-run and
      produced improved answers: task~4 now includes the
      previously-missing TETE$\leftrightarrow$ITRS location
      propagation, and task~5 now includes the prerequisite
      \texttt{self.data.cols = cols} line.  Both moved from
      Near-perfect to Exact, raising the overall score from 160 to 162.

\end{enumerate}
